\documentclass[9pt]{extarticle}

\usepackage[utf8]{inputenc}
\usepackage[a4paper, margin=1.2cm]{geometry}
\usepackage[colorlinks=true, urlcolor=blue]{hyperref}
\usepackage{enumitem}
\usepackage{changepage}

\pagestyle{empty}

\setlist[itemize]{label=\textbullet, noitemsep, topsep=0pt}
\setlength{\parindent}{0pt}

\begin{document}

\begin{center}
	{\huge\bfseries Paul Parisot}
	\\[4pt]
	{\Large Ingénieur Logiciel - C++ \& Optimisation}
	\\
\end{center}

\par\vspace{-5pt}
\rule{\linewidth}{0.4pt}
\par\vspace{-5pt}

\begin{center}
	Github: \href{https://github.com/paulogarithm}{paulogarithm}
	\\
	\href{mailto:paul.parisot@epitech.eu}{paul.parisot@epitech.eu}
	\\
\end{center}

\par\vspace{-10pt}
\rule{\linewidth}{0.4pt}
\par\vspace{7pt}

{\large\bfseries À propos}
\par\vspace{-8pt}
\rule{\linewidth}{0.4pt}
\par\vspace{5pt}

Ingénieur logiciel orienté systèmes, spécialisé en C/C++ (optimisation de performance, programmation réseau, Linux) et en algorithmique. Intérêt marqué pour l'{\bfseries optimisation combinatoire et la recherche opérationnelle}.
\\
Formé à la comptabilité, à l'économie managériale et aux mathématiques à McGill University (Montréal).
\\
Contributeur open-source (plus de 10 000 lignes de code sur des packages critiques). En préparation de l'agrégation d'informatique {\bfseries (algorithmique des graphes, complexité, automates)}.
\par\vspace{5pt}

{\large\bfseries Compétences}
\par\vspace{-8pt}
\rule{\linewidth}{0.4pt}
\par\vspace{5pt}

\hspace*{5pt}{\bfseries Maîtrise}: C/C++, Golang, Python, Lua, Haskell, OCaml, Assembleur, Shell, SQL
\\
\hspace*{5pt}{\bfseries Algorithmique}: graphes, complexité, automates, programmation dynamique, heuristiques
\\
\hspace*{5pt}{\bfseries Outils}: gdb, Valgrind, Git, Docker, CI/CD, Environnement Linux
\\
\hspace*{5pt}{\bfseries Langues}: Français, Anglais (C1, compétences professionnelles)

\par\vspace{5pt}

{\large\bfseries Expérience Professionnelle}
\par\vspace{-8pt}
\rule{\linewidth}{0.4pt}
\par\vspace{5pt}

{\bfseries Gno.land - Ingénieur Logiciel Junior (Open-Source)}
\begin{itemize}[noitemsep, topsep=0pt]
	\item Développement de 6+ packages core en Gnolang (variante de Go pour smart contracts), dont une compatibilité EVM (traduction d'opcodes en instructions Go).
	\item Détection et correction de {\bfseries 6 bugs critiques} causant des crashs runtime ; {\bfseries ~10k lignes de code} contribuées.
	\item Rédaction de tests exhaustifs sous des politiques strictes de CI/CD et de revue de code.
\end{itemize}

\par\vspace{5pt}

{\bfseries Faurecia / Forvia - Ingénieur Démo Logicielle}
\begin{itemize}[noitemsep, topsep=0pt]
	\item Conception d'un {\bfseries moteur de calcul d'itinéraire GPS} et d'une interface de recherche de POI, intégrés à une application démo présentée au {\bfseries CES 2024}.
	\item Correction d'un ancien code {\bfseries C++} pour le faire communiquer avec les bons drivers matériels (LED) ; scripts Python de pilotage de hardware.
	\item Livraison de fonctionnalités prêtes pour la production {\bfseries en avance sur le planning}, salué pour la rapidité et la qualité du code.
\end{itemize}

\par\vspace{5pt}

{\bfseries Epitech / IONIS - Assistant Technique \& Formateur}
\begin{itemize}[noitemsep, topsep=0pt]
	\item Enseignement à {\bfseries plus de 400 étudiants} en C, C++, asm et Haskell (dont des problèmes de multithreading) ; correction et évaluation des rendus.
	\item Création de {\bfseries contenus pédagogiques interactifs} (quiz, cours inversés, tutoriels) et contribution à un {\bfseries taux de réussite historiquement élevé} au sein des promotions coachées en tant qu'assistant technique.
\end{itemize}

\par\vspace{5pt}

{\large\bfseries Projets Personnels}
\par\vspace{-8pt}
\rule{\linewidth}{0.4pt}
\par\vspace{5pt}

\underline{\bfseries \href{https://github.com/paulogarithm/blorox}{BloRox} -- Moteur 3D temps réel en C++}
\par
\begin{adjustwidth}{5pt}{0pt}
	Architecture ECS orientée données : {\bfseries 100k entités animées}, $\sim$6 ms de moteur et $\sim$6 ms de rendu par image. \\
	Alignement mémoire et extensions SIMD du processeur pour la vectorisation ; migration prévue vers {\ttfamily std::simd} (C++26). \\
	Prototypage d'une sérialisation réseau par réflexion statique C++26 (compilateur expérimental).
\end{adjustwidth}
\par
{\bfseries Compétences :} C++23/26, ECS, alignement mémoire, SIMD, optimisation de performance, conception de moteur

\par\vspace{8pt}

\underline{\bfseries Zappy -- Serveur TCP multi-clients en C}
\par
\begin{adjustwidth}{5pt}{0pt}
    Conception et codage seul d'un serveur TCP mono-thread à boucle événementielle ({\bfseries \texttt{select}}), avec gestion des clients par liste chaînée de connexions. \\
    Buffer borné par client et mise en quarantaine temporaire des clients qui le saturent, pour éviter tout épuisement mémoire. \\
    Aucune erreur ni fuite mémoire sous Valgrind (y compris à l'arrêt par SIGINT) ; tests unitaires et fonctionnels.
\end{adjustwidth}
\par
{\bfseries Compétences :} C, sockets TCP, select, gestion mémoire, Valgrind, tests unitaires et fonctionnels

\par\vspace{8pt}

\underline{\bfseries Plazza -- Simulation de pizzerias concurrentes (Epitech)}
\par
\begin{adjustwidth}{5pt}{0pt}
    Simulation en C++ de pizzerias concurrentes reposant sur des threads et des thread pools pour la préparation des commandes.
\end{adjustwidth}

\par\vspace{8pt}

\underline{\bfseries \href{https://github.com/paulogarithm/cclass}{cclass} -- Moteur de mutation mémoire à l'exécution}
\par
\begin{adjustwidth}{5pt}{0pt}
    Framework en C pur réécrivant à l'exécution des instructions assembleur en mémoire ({\bfseries \texttt{mmap}}, flags de protection, alignement strict) pour contourner le surcoût du dispatch par vtable.
\end{adjustwidth}

\par\vspace{8pt}

{\large\bfseries Formation}
\par\vspace{-8pt}
\rule{\linewidth}{0.4pt}
\par\vspace{5pt}

{\bfseries Epitech, Paris} \hfill 2022 -- 2027 \\
Diplôme d'ingénieur en technologies de l'information (diplôme de niveau Master) \\
Certification Niveau 7 - Expert en ingénierie logicielle

\vspace{6pt}

{\bfseries McGill University, Montréal} \hfill 2025 -- 2026 \\
Certificat en Management -- GPA : 3.57/4.0 \\
Cours suivis : Comptabilité, Managerial Economics \& Analysis, Business, Mathématiques, Gestion de projet, Marketing

\end{document}
